% =============================================================================
% Section 7: Limitations
% =============================================================================
\section{Limitations}\label{sec:limitations}

The duplex finding is robust within the {\ISI} framework, but
several structural limitations constrain its external validity.

\subsection{Re-export masking}

Trade-based axes (technology, critical inputs, logistics) use
bilateral import data that record the last country of shipment, not
the country of origin.  Countries that function as re-export hubs
(notably the Netherlands and Belgium) may appear more diversified
than they are, because multiple recorded counterparties ultimately
source from the same upstream producer.  This limitation applies to
Paper~2's full results and is inherited here without modification.

\subsection{Small-economy amplification}

Small economies tend to exhibit higher concentration mechanically:
with fewer trading partners and smaller absolute volumes, the
{\HHI} is structurally elevated.  The top four {\ISI} ranks are
occupied by Malta, Cyprus, Ireland, and Denmark---all small or
medium economies.  The duplex finding (defence and logistics
dominate variance) is partly a consequence of these axes exhibiting
the widest range between small and large economies.

\subsection{Defence data sparsity}

The SIPRI arms transfer database records trend-indicator values for
major conventional weapons.  Countries with no recorded transfers in
the observation window receive imputed scores based on documented
procurement structures.  For countries with minimal arms imports
(e.g., Slovakia, with a defence score of 0.000), the imputed score
may not reflect actual dependence structures.  Slovakia's zero
score is an extreme case that inflates the variance of the defence
axis.

\subsection{Logistics scope}

The logistics axis currently captures containerised freight and port
throughput.  Road, rail, inland waterway, and air freight are
excluded.  For landlocked countries (Austria, Czechia, Hungary,
Slovakia, Luxembourg), the logistics axis may understate
concentration because it omits their primary freight modalities.

\subsection{Static snapshot}

The analysis is based on a single vintage (2024).  The variance
structure may change over time as supply chains evolve, new trade
agreements are signed, or defence procurement shifts.  The duplex
finding should be treated as a snapshot property of the current
data, not as a structural invariant.

\subsection{Equal-weight sensitivity}

The equal-weight assumption ($w_j = 1/6$) is tested by Dirichlet
perturbation (mean $\rho = 0.980$; see \cref{sec:robustness}).
This confirms ranking stability under moderate weight variation,
but does not address the case of deliberately unequal expert
weights, which could reduce or amplify the duplex.  Under a
weighting scheme that down-weights defence and logistics, the
duplex would weaken; under one that up-weights them, it would
intensify.  The finding is conditional on approximate weight parity.
